\chapter{Liouville Conformal Field Theory}
\label{ch:liouville-cft}
\ConventionsEuclideanCFT

\section{Classical Datum and Stress Tensor}

Liouville theory is the noncompact two-dimensional conformal theory whose
interaction is an exponential of a scalar field.  This chapter treats it as a
two-dimensional QFT construction problem whose data are a background surface,
a scalar distribution, a curvature coupling, exponential insertions, Virasoro
Ward identities, degenerate fields, and sewing constraints.  The development
below fixes the normalizations, derives the local representation-theoretic
identities, and states the present theorem boundaries for the probabilistic
and bootstrap constructions.

\begin{definition}[Classical Liouville datum]
\label{def:classical-liouville-datum}
A classical Liouville datum consists of an oriented smooth closed surface
\(\Sigma\), a Riemannian metric \(g\) on \(\Sigma\), a real scalar field
\(\phi\in C^\infty(\Sigma,\mathbb R)\), parameters
\[
  b>0,\qquad Q\in\mathbb R,\qquad \mu>0,
\]
and the Euclidean action
\begin{equation}
  S_L[\phi;g]
  =
  \frac1{4\pi}
  \int_\Sigma \dd^2x\,\sqrt g\,
  \left(
    g^{ab}\partial_a\phi\,\partial_b\phi
    +Q R_g\phi
    +4\pi\mu\,\ee^{2b\phi}
  \right).
  \label{eq:liouville-classical-action}
\end{equation}
Here \(R_g\) is the scalar curvature of \(g\).  The exponential term is a
function of the smooth classical field in this definition; in the quantum
theory it is a renormalized random measure or a normal-ordered local field,
depending on the construction framework.
\end{definition}

Varying a compactly supported smooth \(\delta\phi\) gives
\[
  \delta S_L
  =
  \frac1{4\pi}
  \int_\Sigma\sqrt g\,
  \left(
    -2\Delta_g\phi+Q R_g+8\pi b\mu\ee^{2b\phi}
  \right)\delta\phi .
\]
Thus the classical Euler--Lagrange equation is
\begin{equation}
  -\Delta_g\phi+\frac Q2 R_g+4\pi b\mu\ee^{2b\phi}=0.
  \label{eq:liouville-classical-eom}
\end{equation}
On a flat complex coordinate patch \(z\), the holomorphic stress tensor
representative is
\begin{equation}
  T(z)
  =
  -:\!(\partial\phi)^2\!:(z)+Q\,\partial^2\phi(z).
  \label{eq:liouville-stress-tensor}
\end{equation}
The normal ordering in \eqref{eq:liouville-stress-tensor} is a free-field
short-distance prescription.  It is used here as a perturbative coordinate on
the Virasoro representation data; the nonperturbative probabilistic
construction is stated later as a separate object.

\begin{convention}[Free-field normalization]
\label{conv:liouville-free-field-normalization}
In this chapter the local free field has two-point singularity
\[
  \phi(z)\phi(w)\sim -\frac12\log |z-w|^2,
  \qquad
  \partial\phi(z)\partial\phi(w)\sim -\frac1{2(z-w)^2}.
\]
The holomorphic exponential field is denoted
\[
  V_\alpha(z,\bar z)=\mathopen{:}\ee^{2\alpha\phi(z,\bar z)}\mathclose{:}.
\]
With this convention the standard Liouville relation is
\[
  Q=b+b^{-1}.
  \label{eq:liouville-Q-b-relation}
\]
\end{convention}

\begin{proposition}[Improved central charge and Liouville weights]
\label{prop:liouville-central-charge-weights}
Assume the free-field OPE convention
\ref{conv:liouville-free-field-normalization}.  The stress tensor
\eqref{eq:liouville-stress-tensor} has Virasoro central charge
\begin{equation}
  c_L=1+6Q^2.
  \label{eq:liouville-central-charge}
\end{equation}
The exponential \(V_\alpha\) is Virasoro primary in the free-field
coordinate with conformal weights
\begin{equation}
  h_\alpha=\bar h_\alpha=\alpha(Q-\alpha).
  \label{eq:liouville-exponential-weight}
\end{equation}
\end{proposition}

\begin{proof}
The stress tensor is \(T=T_0+T_Q\), where
\[
  T_0=-:\!(\partial\phi)^2\!:,\qquad T_Q=Q\partial^2\phi.
\]
The double contraction in \(T_0(z)T_0(w)\) gives the central term
\(\frac{1}{2}(z-w)^{-4}\), hence \(c=1\) for \(Q=0\).  The mixed terms have
at most third-order poles and contribute to the transformation of \(T\), not
to the central term.  Since
\[
  \partial^2\phi(z)\partial^2\phi(w)
  =
  \partial_z^2\partial_w^2
  \left(-\frac12\log(z-w)\right)
  \sim
  \frac{3}{(z-w)^4},
\]
the improvement contributes \(3Q^2(z-w)^{-4}\).  Comparing
\[
  T(z)T(w)\sim
  \frac{c/2}{(z-w)^4}
  +\frac{2T(w)}{(z-w)^2}
  +\frac{\partial T(w)}{z-w}
\]
gives \(c=1+6Q^2\).

For the exponential,
\[
  \partial\phi(z)V_\alpha(w,\bar w)
  \sim
  -\frac{\alpha}{z-w}V_\alpha(w,\bar w),
  \qquad
  \partial^2\phi(z)V_\alpha(w,\bar w)
  \sim
  \frac{\alpha}{(z-w)^2}V_\alpha(w,\bar w).
\]
The double contraction of \(-:\!(\partial\phi)^2\!:\) contributes
\(-\alpha^2(z-w)^{-2}V_\alpha\), and the improvement contributes
\(\alpha Q(z-w)^{-2}V_\alpha\).  The simple-pole term is
\((z-w)^{-1}\partial_wV_\alpha\).  Therefore \(V_\alpha\) has the weight
\(\alpha(Q-\alpha)\).
\end{proof}

The interaction \(\mu\ee^{2b\phi}\) is marginal precisely when
\(b(Q-b)=1\), which gives \(Q=b+b^{-1}\).  Thus
\eqref{eq:liouville-Q-b-relation} is an operational normalization condition:
it is the condition that the exponential potential have weight \((1,1)\) in
the background-charge representation.

\section{Exponential Operators and the Probabilistic Object}

Liouville correlation functions are infinite-dimensional Euclidean QFT
objects.  The rigorous probabilistic construction separates the constant mode
from the mean-zero Gaussian free field and interprets the exponential
interaction as Gaussian multiplicative chaos.

\begin{definition}[Seiberg domain in physics normalization]
\label{def:liouville-seiberg-domain}
For insertions \(V_{\alpha_1}(x_1),\ldots,V_{\alpha_n}(x_n)\) on a closed
surface \(\Sigma\) with Euler characteristic \(\chi(\Sigma)\), the Seiberg
domain is the set of momenta satisfying
\begin{equation}
  \alpha_i<\frac Q2\quad\text{for every }i,
  \qquad
  2\sum_{i=1}^n\alpha_i>Q\,\chi(\Sigma).
  \label{eq:liouville-seiberg-bound}
\end{equation}
The first inequalities are local integrability conditions near the marked
points.  The second inequality is the negative-constant-mode integrability
condition.  On the sphere it reads \(\sum_i\alpha_i>Q\).
\end{definition}

The zero-mode check is elementary and fixes the sign conventions.  Write
\(\phi=c+\phi_0\), with \(\int_\Sigma\phi_0\,\mathrm{vol}_g=0\).  The
curvature term contributes \(\exp[-Q\chi(\Sigma)c]\), while the insertions
contribute \(\exp[2c\sum_i\alpha_i]\).  As \(c\to-\infty\), integrability
requires \(2\sum_i\alpha_i-Q\chi(\Sigma)>0\).  The positive \(c\) direction
is controlled by the Liouville potential when its multiplicative-chaos
measure is nonzero.

\begin{definition}[Gaussian multiplicative chaos input]
\label{def:liouville-gmc-input}
Let \(X_g\) be a mean-zero Gaussian free field on \((\Sigma,g)\), normalized
by a Green kernel with logarithmic singularity \(-\log d_g(x,y)\).  For
\(\gamma\in(0,2)\), let \(X_{g,\varepsilon}\) be a smooth mollification at
scale \(\varepsilon\).  The subcritical Gaussian multiplicative-chaos measure
is the probability-limit object
\begin{equation}
  M_{\gamma,g}(\dd x)
  =
  \lim_{\varepsilon\downarrow0}
  \varepsilon^{\gamma^2/2}
  \ee^{\gamma X_{g,\varepsilon}(x)}\,\mathrm{vol}_g(\dd x)
  \label{eq:liouville-gmc-measure}
\end{equation}
in the topology of weak convergence in probability against continuous test
functions.  In Liouville theory \(\gamma=2b\) after translating from the
physics field \(\phi\) of Convention
\ref{conv:liouville-free-field-normalization} to the probabilistic GFF
normalization.
\end{definition}

\begin{quotedtheorem}[Probabilistic Liouville correlation functions]
\label{qthm:gkrv-liouville-correlators}
Guillarmou--Kupiainen--Rhodes--Vargas construct Liouville correlation
functions on compact Riemann surfaces in the subcritical GMC regime by
integrating the constant mode against the Gaussian multiplicative-chaos
functional, with insertions in the Seiberg domain
\eqref{eq:liouville-seiberg-bound}.  The construction gives finite
correlation functions, conformal covariance with central charge
\(1+6Q^2\), and meromorphic continuation in the insertion momenta after the
initial probabilistic domain is established.  This chapter uses the theorem
as the measure-theoretic construction of the Liouville path integral and then
derives the Virasoro consequences needed below.
\end{quotedtheorem}

\begin{definition}[Reflection relation]
\label{def:liouville-reflection-relation}
The Liouville spectrum is parametrized by
\[
  \alpha=\frac Q2+\ii P,\qquad P\in\mathbb R_{\ge0},
\]
with conformal weight \(Q^2/4+P^2\).  The reflection relation is the
identification, at the level of meromorphic correlation functions,
\begin{equation}
  V_\alpha=R_b(\alpha)\,V_{Q-\alpha},
  \label{eq:liouville-reflection-relation}
\end{equation}
where \(R_b(\alpha)\) is the meromorphic reflection amplitude fixed by the
two-point normalization and equivalently by the DOZZ three-point function.
Thus a single bulk primary label is represented by the reflection orbit
\(\{\alpha,Q-\alpha\}\), with the normalization encoded in \(R_b(\alpha)\).
\end{definition}

\section{The DOZZ Formula and Shift Checks}

The exact three-point function is the Dorn--Otto--Zamolodchikov--Zamolodchikov
structure constant.  It is theorem-level input for this chapter, but the
special functions and convention checks are part of the manuscript.

\begin{definition}[The \(\Upsilon_b\) function]
\label{def:upsilon-b-function}
For \(b>0\) and \(Q=b+b^{-1}\), \(\Upsilon_b(x)\) is the entire function
characterized by self-duality \(\Upsilon_b(x)=\Upsilon_{1/b}(x)\), reflection
\(\Upsilon_b(x)=\Upsilon_b(Q-x)\), zeros at
\[
  x=-mb-nb^{-1}
  \quad\text{and}\quad
  x=Q+mb+nb^{-1},
  \qquad m,n\in\mathbb Z_{\ge0},
\]
and shift relations
\begin{align}
  \Upsilon_b(x+b)
  &=
  \gamma(bx)\,b^{1-2bx}\,\Upsilon_b(x),
  \label{eq:upsilon-b-shift-b}\\
  \Upsilon_b(x+b^{-1})
  &=
  \gamma(x/b)\,b^{2x/b-1}\,\Upsilon_b(x),
  \label{eq:upsilon-b-shift-binverse}
\end{align}
where \(\gamma(x)=\Gamma(x)/\Gamma(1-x)\).  These equations fix the
normalization up to the conventional value of \(\Upsilon'_b(0)\), which is
kept explicit in the three-point constant.
\end{definition}

\begin{quotedtheorem}[DOZZ three-point function]
\label{qthm:dozz-formula}
In the normalization of Convention
\ref{conv:liouville-free-field-normalization}, the sphere three-point
function has the form
\[
  \left\langle
    V_{\alpha_1}(z_1)V_{\alpha_2}(z_2)V_{\alpha_3}(z_3)
  \right\rangle_{\mathbb P^1}
  =
  \frac{
    C_b(\alpha_1,\alpha_2,\alpha_3)
  }{
    |z_{12}|^{2(h_1+h_2-h_3)}
    |z_{13}|^{2(h_1+h_3-h_2)}
    |z_{23}|^{2(h_2+h_3-h_1)}
  },
\]
with \(h_i=\alpha_i(Q-\alpha_i)\) and
\begin{equation}
\begin{aligned}
  C_b(\alpha_1,\alpha_2,\alpha_3)
  &=
  \left[
    \pi\mu\,\gamma(b^2)\,b^{2-2b^2}
  \right]^{(Q-\alpha_1-\alpha_2-\alpha_3)/b}
  \\
  &\quad\times
  \frac{
    \Upsilon'_b(0)
    \prod_{i=1}^3\Upsilon_b(2\alpha_i)
  }{
    \Upsilon_b(\alpha_1+\alpha_2+\alpha_3-Q)
    \Upsilon_b(\alpha_1+\alpha_2-\alpha_3)
    \Upsilon_b(\alpha_1+\alpha_3-\alpha_2)
    \Upsilon_b(\alpha_2+\alpha_3-\alpha_1)
  } .
\end{aligned}
  \label{eq:dozz-structure-constant}
\end{equation}
The formula was proposed by Dorn--Otto and by the Zamolodchikovs and was
proved in the probabilistic construction by Kupiainen--Rhodes--Vargas.  This
chapter quotes the theorem rather than reproducing the GMC proof.
\end{quotedtheorem}

The shift relations \eqref{eq:upsilon-b-shift-b}--\eqref{eq:upsilon-b-shift-binverse}
are the convention-sensitive part of \eqref{eq:dozz-structure-constant}.  They
are what make the degenerate-field shift equations compatible with the
normalization of the exponential interaction.  The companion script
\texttt{calculation-checks/liouville\_bpz\_checks.py} verifies the level-two
null-vector arithmetic used below; a later check should expand this to
symbolic DOZZ shift-ratio simplifications.

\section{Degenerate Insertions and BPZ Equations}

The simplest degenerate Liouville momentum is \(\alpha=-b/2\).  Its conformal
weight is
\[
  h_{-b/2}
  =
  -\frac12-\frac{3b^2}{4}.
\]

\begin{proposition}[Level-two Liouville null vector]
\label{prop:liouville-level-two-null-vector}
Let \(c=1+6Q^2\), \(Q=b+b^{-1}\), and
\[
  h=-\frac12-\frac{3b^2}{4}.
\]
In a highest-weight Virasoro module with highest vector \(|h\rangle\), the
level-two vector
\begin{equation}
  \left(L_{-1}^2+b^2L_{-2}\right)|h\rangle
  \label{eq:liouville-level-two-null-vector}
\end{equation}
is annihilated by \(L_1\) and \(L_2\).  Hence it is a null vector whenever the
module is quotiented to the degenerate representation carried by
\(V_{-b/2}\).
\end{proposition}

\begin{proof}
For a candidate \((L_{-1}^2+\kappa L_{-2})|h\rangle\), the Virasoro
commutators give
\[
  L_1L_{-1}^2|h\rangle=(4h+2)L_{-1}|h\rangle,
  \qquad
  L_1L_{-2}|h\rangle=3L_{-1}|h\rangle,
\]
and
\[
  L_2L_{-1}^2|h\rangle=6h|h\rangle,
  \qquad
  L_2L_{-2}|h\rangle=(4h+c/2)|h\rangle .
\]
The \(L_1\) equation requires
\(\kappa=-(4h+2)/3=b^2\).  Substituting
\(h=-1/2-3b^2/4\) and \(c=1+6(b+b^{-1})^2\) into the \(L_2\) equation gives
\[
  6h+b^2(4h+c/2)=0.
\]
Thus both positive modes annihilate \eqref{eq:liouville-level-two-null-vector}.
\end{proof}

\begin{proposition}[BPZ differential equation]
\label{prop:liouville-bpz-equation}
Let
\[
  F(z;z_1,\ldots,z_n)
  =
  \left\langle
    V_{-b/2}(z)\prod_{i=1}^n V_{\alpha_i}(z_i)
  \right\rangle
\]
be a separated-point correlation function in a domain where the Virasoro Ward
identity and the degenerate null-vector quotient are valid.  Then \(F\)
satisfies
\begin{equation}
  \left[
    \partial_z^2
    +b^2\sum_{i=1}^n
    \left(
      \frac{h_i}{(z-z_i)^2}
      +\frac1{z-z_i}\partial_{z_i}
    \right)
  \right]F=0,
  \qquad
  h_i=\alpha_i(Q-\alpha_i).
  \label{eq:liouville-bpz-equation}
\end{equation}
\end{proposition}

\begin{proof}
The \(L_{-1}\) insertion at \(z\) is \(\partial_z\).  The Ward identity for a
stress-tensor contour around \(z\), moved to contours around the other marked
points, gives
\[
  \left\langle
    (L_{-2}V_{-b/2})(z)\prod_iV_{\alpha_i}(z_i)
  \right\rangle
  =
  \sum_i
  \left(
    \frac{h_i}{(z-z_i)^2}
    +\frac1{z-z_i}\partial_{z_i}
  \right)F.
\]
Inserting the null-vector equation
\((L_{-1}^2+b^2L_{-2})V_{-b/2}=0\) gives
\eqref{eq:liouville-bpz-equation}.  This is an identity for
separated-point correlators; coincident-point limits require the OPE and
contact-term renormalization.
\end{proof}

The BPZ equation constrains correlators recursively because fusing
\(V_{-b/2}\) with a generic \(V_\alpha\) has only two allowed Liouville
momenta,
\[
  V_{-b/2}V_\alpha
  \sim
  C_+(\alpha)V_{\alpha-b/2}
  +
  C_-(\alpha)V_{\alpha+b/2}.
\]
Solving \eqref{eq:liouville-bpz-equation} near one puncture and matching it
to the same equation near another puncture gives shift equations for the
three-point structure constants.  The DOZZ constant is the meromorphic
solution with the reflection, shift, and analyticity properties stated above.

\section{Sewing Status and Open Boundary}

The closed-surface probabilistic theory, the bootstrap formula, and the
Segal/Kontsevich--Segal sewing formulation require explicit comparison maps.

\begin{openproblem}[Liouville sewing and functorial closure]
\label{op:liouville-functorial-closure-volume-v}
Construct, in a single framework, the Hilbert or topological vector spaces
assigned to circles, the multilinear maps assigned to bordered Riemann
surfaces, the boundary state spaces, and the sewing maps whose closed-surface
matrix elements reproduce the probabilistic Liouville correlators and the
DOZZ/bootstrap conformal blocks.  The comparison must specify the topology on
continuous-spectrum direct integrals, the domain of unbounded vertex
operators, and the convergence of the gluing integrals over
\(\alpha=Q/2+\ii P\).  This is the Volume V form of
the Liouville and nonrational-CFT open problem in the
Kontsevich--Segal functorial-QFT chapter.
\end{openproblem}

Thus the status is deliberately layered.  The classical action and
background-charge Ward algebra are derived in the chapter.  The
measure-theoretic path integral and the DOZZ formula are quoted theorem
inputs with named proof boundaries.  Full functorial sewing remains an open
problem for the monograph's foundational program.
